% Системные_требования.tex — Системные требования (DOC-05), версия 0.1 — СКЕЛЕТ
% Комплект документации криптотренажёра mephi-12 (НИЯУ МИФИ, курсовая работа)
% Стиль — ДО-178C/КТ-178C; преамбула и титул — по единому шаблону full_doc_template.tex
% НАЗНАЧЕНИЕ v0.1: зафиксировать структуру, методику ID и шаблон карточки требования,
% чтобы 6 авторов в Спринте 2 (задачи 2.6–2.8 плана) наполняли документ единообразно.
% Наполнение разделов 4 и 5 требованиями — Спринт 2; v1.0 — до 05.11.2026.
\documentclass[12pt,a4paper]{article}

% =============================================================================
% КОНФУГУРАЦИЯ И ПРЕАМБУЛА (строго по единому шаблону комплекта)
% =============================================================================

% --- конфигурация метаданных ---
% ---------- Прямые значения ----------
\newcommand{\Company}{НИЯУ МИФИ}
\newcommand{\DocumentTitleMain}{Системные требования}
\newcommand{\DocumentTitleSub}{для платформы изучения и практики курса по криптографии в МИФИ (криптотренажёр mephi-12)}
\newcommand{\Version}{0.1}
\newcommand{\DocDate}{24.10.2026}
\newcommand{\AuthorName}{Анохин Г.А.}
\newcommand{\AgreedBy}{}
\newcommand{\ApprovedBy}{}
\newcommand{\City}{Москва}

% --- преамбула и стили ---
\usepackage[utf8]{inputenc}
\usepackage[T2A]{fontenc}
\usepackage[russian]{babel}

% Times New Roman-like
\usepackage{newtxtext,newtxmath} % похож на Times; обеспечивает текст+математику
\usepackage{setspace}
\onehalfspacing % межстрочный интервал 1.5

% Страницы
\usepackage{geometry}
\geometry{left=25mm,right=15mm,top=20mm,bottom=20mm}

% отступ первой строки
\usepackage{indentfirst}
\setlength{\parindent}{1.25cm}

% таблицы, картинки, списки, подписи
\usepackage{graphicx}
\usepackage{caption}
\usepackage{subcaption}
\captionsetup[table]{position=top}    % подпись таблицы над ней
\captionsetup[figure]{position=bottom}% подпись рисунка под ней
\usepackage{booktabs}
\usepackage{tabularx}
\usepackage{longtable}
\usepackage{array}
\usepackage{enumitem}
\setlist[itemize]{left=0pt,label=\textbullet}
\setlist[enumerate]{left=0pt,label=\arabic*.}

% код и конфигурации
\usepackage{listings}
\lstset{basicstyle=\ttfamily\small,breaklines=true}

% гиперссылки
\usepackage{hyperref}
\hypersetup{colorlinks=true,linkcolor=black,urlcolor=blue,
  pdftitle={Системные требования — криптотренажёр mephi-12},
  pdfauthor={Анохин Г.А.}}

\pagestyle{plain}  % только номер страницы внизу по центру

% нумерация разделов — стандартная десятичная (article по умолчанию)
% дополнительные удобства
\usepackage{tocloft} % управление оглавлением
\setlength{\cftbeforesecskip}{6pt}

% небольшие утилиты
\newcommand{\reqfield}[1]{\textbf{#1}:} % для шаблонов в приложениях

% Карточка требования: двухколоночная tabularx «поле — содержание» (раздел 3).
% Первый столбец — имя поля полужирным, второй — содержание с переносами.
% ВАЖНО: tabularx нельзя прятать внутрь пользовательского окружения (сканирует
% литеральное \end{tabularx}), а пакет array не раскрывает макросы в спецификации
% столбцов — поэтому карточки оформляются явным окружением tabularx с литальной
% спецификацией столбцов \begin{center}\small\begin{tabularx}{\textwidth}{...} ... \end{tabularx}\end{center}.

\begin{document}

% =============================================================================
% ТИТУЛЬНЫЙ ЛИСТ
% =============================================================================
\begin{titlepage}
  \thispagestyle{empty} % убирает номер/колонтитул на этой странице
  \centering
  \vspace*{5.8cm}
  {\Huge \textbf{\DocumentTitleMain}\par}
  {\Large \vspace{4pt} \DocumentTitleSub \par} % вторая часть — обычным шрифтом, чуть ниже
  \vspace{0.8cm}

  {\large Версия: \Version \par}
  {\large Дата: \DocDate \par}
  \vspace{1.2cm}

  % ----- Таблица с ролями -----
  \centering
  \begin{tabularx}{0.95\textwidth}{@{\hspace{6pt}}p{5.5cm} p{6cm} p{4cm}@{\hspace{6pt}}}
    \toprule
    \textbf{Роль} & \textbf{ФИО} & \textbf{Подпись / Дата} \\
    \midrule
    Разработал: & \AuthorName & \\
    \addlinespace[6pt]
    Согласовано: & \AgreedBy & \\
    \addlinespace[6pt]
    Утверждено: & \ApprovedBy & \\
    \bottomrule
  \end{tabularx}

  \vfill
  {\large г. \City, \the\year \par}
\end{titlepage}

% =============================================================================
% ЛИСТ РЕГИСТРАЦИИ ИЗМЕНЕНИЙ
% =============================================================================
\section*{Лист регистрации изменений}
\addcontentsline{toc}{section}{Лист регистрации изменений}

\begin{table}[h!]
\centering
\small
\begin{tabularx}{\textwidth}{
  >{\centering\arraybackslash}p{1.1cm}
  >{\centering\arraybackslash}p{2.1cm}
  p{2.3cm}
  p{2.6cm}
  p{1.6cm}
  X
  p{1.8cm}
}
\toprule
\textbf{Версия} & \textbf{Дата} & \textbf{Автор} & \textbf{Изменённый элемент} & \textbf{Х-р изменения} & \textbf{Описание изменения} & \textbf{Утвердил} \\
\midrule
0.1 & 24.10.2026 & Анохин Г.А. & Весь документ & Добавлено & Скелет документа (Спринт 1): структура, методика ID, шаблон карточки требования. Наполнение — Спринт 2, v1.0 — до 05.11.2026. & \\
\bottomrule
\end{tabularx}
\end{table}

\clearpage
\tableofcontents
\clearpage

% =============================================================================
% ОСНОВНОЕ СОДЕРЖАНИЕ
% =============================================================================

\section{Введение \texttt{[SR.INTRO]}}\label{sec:intro}

\subsection{Назначение документа \texttt{[SR.INTRO.PURPOSE]}}
Настоящий документ содержит системные требования (далее по тексту — СТ, требования SR-XXX) к
системе «Платформа для изучения и практики курса по криптографии в МИФИ» (краткое наименование —
криптотренажёр mephi-12; далее по тексту — Система). Документ декомпозирует положения
технического задания DOC-02 (маркеры \texttt{[TS.*]}) в атомарные проверяемые требования и
является для Системы документом уровня системных требований в смысле ДО-178C/КТ-178C
(стандарты программных требований применяются к следующему уровню — DOC-06).

Версия 0.1 представляет собой скелет документа: зафиксированы структура разделов, методика
идентификации требований, шаблон карточки требования и правила её заполнения, а также три
карточки-образца. Наполнение разделов~\ref{sec:func} и~\ref{sec:nft} требованиями выполняется
владельцами доменов в Спринте~2 (до 03.11--04.11.2026), редакторская сборка v1.0 — до
05.11.2026 (раздел~\ref{sec:own}).

\subsection{Область применения \texttt{[SR.INTRO.SCOPE]}}
Документ охватывает всю функциональность Системы: аутентификацию и пользователей, сессии задач
и планировщик автопровала, ранцевые криптосистемы (супервозрастающая и на степенях), бэкдор в
ранцевой криптосистеме, схему Эль-Гамаля с математической библиотекой, конфигурации задач и
модерацию, а также нефункциональные требования: безопасность, интеграцию и интерфейсы,
надёжность и эксплуатацию, производительность.

Объект требований — монолит Kotlin/Spring Boot: REST API с базовым путём \texttt{/api} на
порту 8083 (7 контроллеров, 14 эндпоинтов), 8 JPA-сущностей, роли USER/ADMIN,
аутентификация JWT, СУБД PostgreSQL~17, развёртывание средствами Docker. Типы задач сессий:
\begin{itemize}
  \item \texttt{BACKPACK\_CODE\_SUPER\_INCREASING} — ранцевая КС на супервозрастающей
        последовательности;
  \item \texttt{BACKPACK\_CODE\_DEGREES} — ранцевая КС на степенях;
  \item \texttt{BACKPACK\_BACKDOOR} — бэкдор в ранцевой КС;
  \item \texttt{EL\_GAMAL} — схема Эль-Гамаля.
\end{itemize}
Лимит времени на задачу — 20~минут; автоматический провал просроченных задач выполняется
планировщиком \texttt{ProblemAutoFailScheduler} с интервалом 60~с.

Вне области настоящего документа: детали реализации (структуры БД, классы, компоненты ПО) —
уровень программных требований DOC-06; тестовый код и протоколы прогонов (исключены из
комплекта решением от 20.10.2026); управление конфигурацией документа — DOC-03.

\subsection{Связь с другими документами \texttt{[SR.INTRO.DOCS]}}
Входом для настоящего документа является техническое задание DOC-02 (маркеры \texttt{[TS.*]}
указываются в поле «Источник» каждой карточки требования). Выходами — программные требования
DOC-06/DOC-06.1--06.6, матрица трассировки DOC-09 и верификационные кейсы DOC-10
(таблица~\ref{tab:docs}).

\begin{table}[h!]
\centering
\caption{Связь настоящего документа с комплектом документации}
\label{tab:docs}
\begin{tabularx}{\textwidth}{p{2.2cm}>{\raggedright\arraybackslash}p{6.2cm}>{\raggedright\arraybackslash}X}
\toprule
\textbf{Документ} & \textbf{Название} & \textbf{Роль по отношению к настоящим СТ} \\
\midrule
DOC-02 & Техническое задание, \texttt{[TS.*]} & Вход: первоисточник; каждое SR ссылается на маркер \texttt{[TS.*]} в поле «Источник» \\
DOC-06 & Программные требования (сводно) & Выход: детализация SR в SWR (DOC-06.1--06.6 — по подсистемам) \\
DOC-09 & Матрица трассировки & Выход: цепочка SR~$\to$~SWR~$\to$~TC (CSV-источник) \\
DOC-10 & План верификации ПО & Выход: поле «Сценарий проверки» карточки — основа кейсов TC-XXX \\
DOC-03 & ПУКПО & Управление версиями и базовыми версиями настоящего документа \\
DOC-04 & Процедура формальной инспекции & Контроль качества: ФИ-СТ (06--07.11.2026), протокол DOC-13.3 \\
\bottomrule
\end{tabularx}
\end{table}

\subsection{Термины и сокращения \texttt{[SR.INTRO.TERMS]}}
Полные определения всех терминов Системы (Сессия задач, Задача, Тип задачи, Планировщик
автопровала, Конфигурация задач, JWT и др.) поддерживаются в сквозном глоссарии комплекта —
рабочем материале \texttt{notes/methodology/id-methodology.md} (владелец — Анохин Г.А.);
при расхождении определений приоритет имеет глоссарий комплекта, а расхождение является
предметом замечания ФИ (категория «непротиворечивость»). Ниже — только сокращения,
используемые в настоящем документе.

\begin{description}
  \item[SR:] системное требование; идентификаторы SR-XXX, маркеры разделов \texttt{[SR.*]} (настоящий документ, DOC-05).
  \item[SWR:] программное требование; идентификаторы SWR-XXX, маркеры \texttt{[SWR.*]} (DOC-06, DOC-06.1--06.6).
  \item[TC:] верификационный кейс TC-XXX из плана верификации ПО (DOC-10).
  \item[DEF:] замечание формальной инспекции DEF-XXX (реестр \texttt{defects.md}).
  \item[ФИ:] формальная инспекция по методу Фагана (процедура DOC-04, протоколы DOC-13.1--13.6).
  \item[СТ:] системные требования (синоним SR; настоящий документ).
  \item[ТЗ:] техническое задание (DOC-02).
  \item[НФТ:] нефункциональные требования (раздел~\ref{sec:nft} настоящего документа).
\end{description}

\section{Методика идентификации требований \texttt{[SR.METH]}}\label{sec:meth}
Настоящий раздел сводит действующую методику комплекта
\texttt{notes/methodology/id-methodology.md} (v1.0 от 24.10.2026) к правилам, обязательным при
наполнении настоящего документа; методика обязательна для всех авторов, при расхождении —
применяется методика.

\subsection{Схема идентификаторов \texttt{[SR.METH.ID]}}
Идентификаторы требований — плоские сквозные с ведущими нулями: \texttt{SR-001}\ldots\texttt{SR-999}
(при исчерпании диапазона расширение до четырёх знаков оформляет редактор СТ единым решением).
Доменные префиксы в самих ID не используются: принадлежность требованию домену выражается
разделом документа и маркером \texttt{[SR.*]}, а не номером. Сквозную уникальность идентификаторов
по всему DOC-05 ведёт редактор СТ (Анохин Г.А.); реестр требований ведётся в приложении к
настоящему документу и дублируется в матрице DOC-09.

\subsection{Правила присвоения и жизненный цикл идентификатора \texttt{[SR.METH.LIFE]}}
\begin{enumerate}
  \item \textbf{Уникальность навсегда.} ID уникален по всему документу и не выдаётся повторно
        ни при каких условиях.
  \item \textbf{Удалённое не переиспользуется.} Удалённое требование остаётся в реестре и в
        матрице DOC-09 с пометкой «удалено в v\ldots» и указанием причины; дыры в нумерации
        допустимы и не заполняются.
  \item \textbf{Изменение смысла = новый ID.} Если правка меняет проверяемое утверждение
        (объект, условие, меру), заводится новый ID, старый закрывается по правилу~2; переход
        фиксируется в листе регистрации изменений и в DOC-09. Редакционная правка (стиль,
        опечатки, уточнение без изменения смысла) сохраняет ID.
  \item \textbf{Присвоение до публикации.} ID присваивает владелец домена при включении
        требования в черновик своего раздела; редактор контролирует уникальность и
        резервирование и самостоятельно требования не нумерует (кроме своих доменов).
\end{enumerate}

\subsection{Маркеры разделов \texttt{[SR.METH.MARKERS]}}
Маркер ставится в заголовке раздела и однозначно указывает домен: \texttt{[SR.РАЗДЕЛ]} —
латиница, верхний регистр, разделитель «точка», глубина не более трёх уровней (возможно
уточнение третьим уровнем, например \texttt{[SR.AUTH.JWT]}). Маркер относится к разделу
целиком, включая все его карточки требований. Перечень маркеров настоящего документа —
таблица~\ref{tab:markers}; новые доменные коды вводит только редактор СТ с фиксацией в
методике.

\begin{table}[h!]
\centering
\caption{Доменные маркеры DOC-05 и владельцы разделов}
\label{tab:markers}
\begin{tabularx}{\textwidth}{p{3.3cm}>{\raggedright\arraybackslash}X p{2.6cm}}
\toprule
\textbf{Маркер} & \textbf{Раздел системных требований} & \textbf{Автор раздела} \\
\midrule
\texttt{[SR.AUTH]} & Регистрация, аутентификация, JWT, роли, профиль (п.~\ref{sec:auth}) & Кокорев В.Д. \\
\texttt{[SR.SESSIONS]} & Сессии задач: создание, лимиты, проверка решений, автопровал (п.~\ref{sec:sessions}) & Анохин Г.А. \\
\texttt{[SR.BACKPACK]} & Ранцевые КС: супервозрастающая и на степенях (п.~\ref{sec:backpack}) & Георгиев Д.Д. \\
\texttt{[SR.BACKDOOR]} & Бэкдор в ранцевой КС (п.~\ref{sec:backdoor}) & Минасов М.Р. \\
\texttt{[SR.ELGAMAL]} & Схема Эль-Гамаля и матбиблиотека (п.~\ref{sec:elgamal}) & Портнов Л.М. \\
\texttt{[SR.CONFIG]} & Конфигурации задач, группы, модерация (п.~\ref{sec:config}) & Самбьюи А.Ж. \\
\texttt{[SR.SEC]} & НФТ: безопасность (п.~\ref{sec:sec}) & Кокорев В.Д. \\
\texttt{[SR.INT]} & НФТ: интеграция, интерфейсы (п.~\ref{sec:int}) & Анохин Г.А. \\
\texttt{[SR.REL]} & НФТ: надёжность, эксплуатация, развёртывание (п.~\ref{sec:rel}) & Самбьюи А.Ж. \\
\texttt{[SR.PERF]} & НФТ: производительность (п.~\ref{sec:perf}) & по решению редактора СТ \\
\bottomrule
\end{tabularx}
\end{table}

\subsection{Правила формулировки требований \texttt{[SR.METH.WORDING]}}
\begin{enumerate}
  \item \textbf{Атомарность.} Одно требование — одна атомарная проверяемая формулировка;
        требование проверяется одним кейсом TC одного типа (позитивный/негативный) без
        разделения на случаи.
  \item \textbf{Каноническая форма.} Функциональное требование: «Система должна
        \textless{}действие\textgreater{} при \textless{}условие\textgreater{}»; для НФТ — с
        измеримой мерой («в течение не более $N$\ldots», «не менее $N$\ldots»).
  \item \textbf{Запрещены связки «и/или»} и составные требования внутри одного ID:
        «Система должна A и B» оформляется двумя требованиями с отдельными ID и TC.
  \item \textbf{Проверяемость.} Формулировка допускает объективную проверку по сценарию из
        карточки; термины — только из глоссария комплекта, в единственном написании;
        эндпоинты, состояния и типы — в точности как в коде Системы.
  \item \textbf{Без деталей реализации.} SR не содержит структур БД, имён классов и бинов —
        это уровень SWR (DOC-06); SWR, в свою очередь, не вводит функций, отсутствующих в SR.
\end{enumerate}

\section{Структура требований и шаблон карточки \texttt{[SR.CARD]}}\label{sec:card}
Каждое требование настоящего документа оформляется карточкой по единому шаблону
(п.~\ref{sec:cardtmpl}); правила заполнения полей — п.~\ref{sec:cardrules}; образцы
заполнения — п.~\ref{sec:cardex}.

\subsection{Шаблон карточки требования \texttt{[SR.CARD.TEMPLATE]}}\label{sec:cardtmpl}
Карточка представляет собой таблицу «поле — содержание» из девяти полей. Пустой шаблон
приведён ниже; при наполнении разделов~\ref{sec:func} и~\ref{sec:nft} каждая карточка
оформляется копией этого шаблона.

\begin{center}
\small
\begin{tabularx}{\textwidth}{>{\bfseries}p{3.4cm}>{\raggedright\arraybackslash}X}
\toprule
ID & SR-XXX \\
\addlinespace
Название & (краткое имя существа требования, до 5 слов) \\
\addlinespace
Источник & \texttt{[TS.\ldots]} (маркер раздела DOC-02) \\
\addlinespace
Приоритет & О / Ж \\
\addlinespace
Формулировка & Система должна \ldots{} при \ldots \\
\addlinespace
Сценарий проверки & (способ проверки с ожидаемым результатом; основа TC) \\
\addlinespace
Зависимости & SR-XXX, \ldots{} (или «---») \\
\addlinespace
Владелец & \texttt{[SR.\ldots]}, Фамилия И.О. \\
\addlinespace
Статус & Черновик \\
\bottomrule
\end{tabularx}
\end{center}

\subsection{Правила заполнения полей \texttt{[SR.CARD.RULES]}}\label{sec:cardrules}
\begin{itemize}
  \item \reqfield{ID} плоский сквозной номер по разделу~\ref{sec:meth}; присваивает владелец
        домена, сквозную уникальность контролирует редактор.
  \item \reqfield{Название} краткое уникальное имя требования в инфинитивной форме.
  \item \reqfield{Источник} маркер \texttt{[TS.*]} раздела DOC-02, из положения которого
        декомпозируется требование; при нескольких источниках перечисляются через запятую.
  \item \reqfield{Приоритет} «О» — обязательное для приёмки; «Ж» — желательное (реализация
        по возможности, невыполнение не блокирует приёмку).
  \item \reqfield{Формулировка} каноническая форма «Система должна \ldots» по
        п.~\ref{sec:meth}; атомарная, проверяемая, без деталей реализации.
  \item \reqfield{Сценарий проверки} краткое описание способа проверки с ожидаемым результатом
        (запрос/действие $\to$ наблюдаемый исход); основа верификационного кейса TC в DOC-10.
  \item \reqfield{Зависимости} другие SR-XXX, без которых требование не реализуемо или не
        проверяемо; при отсутствии — прочерк.
  \item \reqfield{Владелец} маркер домена \texttt{[SR.*]} и ФИО владельца домена из
        таблицы~\ref{tab:markers} (не обязательно автор карточки).
  \item \reqfield{Статус} Черновик (поставил владелец) / Согласовано (редактор СТ включил в
        сборку) / Утверждено (ФИ-СТ вынесла решение).
\end{itemize}

\subsection{Примеры-образцы заполнения \texttt{[SR.CARD.EXAMPLES]}}\label{sec:cardex}
Ниже приведены три полностью заполненные карточки-образца (по одной из разделов
\texttt{[SR.AUTH]}, \texttt{[SR.SESSIONS]}, \texttt{[SR.BACKPACK]}), демонстрирующие форму и
глубину проработки. Идентификаторы SR-001--SR-003 закреплены за этими требованиями и в
Спринте~2 не переносятся на другие. Обратите внимание на демонстрацию правила атомарности
(п.~\ref{sec:meth}): смежные аспекты каждой функции (например, выдача JWT при регистрации —
\texttt{[SR.AUTH]}, хранение пароля в виде BCrypt-хеша — \texttt{[SR.SEC]}) оформляются
отдельными карточками.

\emph{(пример-образец, финальная редакция — Спринт 2)}

\begin{center}
\small
\begin{tabularx}{\textwidth}{>{\bfseries}p{3.4cm}>{\raggedright\arraybackslash}X}
\toprule
ID & SR-001 \\
\addlinespace
Название & Уникальность учётной записи \\
\addlinespace
Источник & \texttt{[TS.REQ.FUNC.AUTH.REG]} (DOC-02) \\
\addlinespace
Приоритет & О — обязательное \\
\addlinespace
Формулировка & Система должна отклонять регистрацию с уже используемым \texttt{email},
возвращая ошибку валидации. \\
\addlinespace
Сценарий проверки & Повторная отправка \texttt{POST /api/auth/registration} с уже
зарегистрированным \texttt{email} --- ожидается отказ (ошибка валидации), новая учётная
запись не создаётся (негативный кейс TC в DOC-10). Позитивный сценарий регистрации ---
отдельное требование (Спринт~2). \\
\addlinespace
Зависимости & --- \\
\addlinespace
Владелец & \texttt{[SR.AUTH]}, Кокорев В.Д. \\
\addlinespace
Статус & Черновик \\
\bottomrule
\end{tabularx}
\end{center}

\emph{(пример-образец, финальная редакция — Спринт 2)}

\begin{center}
\small
\begin{tabularx}{\textwidth}{>{\bfseries}p{3.4cm}>{\raggedright\arraybackslash}X}
\toprule
ID & SR-002 \\
\addlinespace
Название & Автоматический провал просроченной задачи \\
\addlinespace
Источник & \texttt{[TS.REQ.FUNC.SESSION.TIME]} (DOC-02) \\
\addlinespace
Приоритет & О --- обязательное \\
\addlinespace
Формулировка & Система должна автоматически переводить задачу сессии из состояния
\texttt{NEW} в состояние \texttt{FAILED}, если с момента создания задачи прошло более
20~минут. \\
\addlinespace
Сценарий проверки & Создать сессию хотя бы с одной задачей и не отправлять решение; после
истечения 20~минут с момента \texttt{created\_at} задачи убедиться, что очередной цикл
планировщика (интервал 60~с) перевёл задачу в \texttt{FAILED}, а сессия без задач в
\texttt{NEW} закрыта. Проверка --- по данным БД и ответу
\texttt{GET /api/moderator/tasks/session}. \\
\addlinespace
Зависимости & SR-0XX «Создание сессии по конфигурации группы» (назначается в Спринте~2,
\texttt{[SR.SESSIONS]}). \\
\addlinespace
Владелец & \texttt{[SR.SESSIONS]}, Анохин Г.А. \\
\addlinespace
Статус & Черновик \\
\bottomrule
\end{tabularx}
\end{center}

\emph{(пример-образец, финальная редакция — Спринт 2)}

\begin{center}
\small
\begin{tabularx}{\textwidth}{>{\bfseries}p{3.4cm}>{\raggedright\arraybackslash}X}
\toprule
ID & SR-003 \\
\addlinespace
Название & Генерация редакционной ранцевой задачи \\
\addlinespace
Источник & \texttt{[TS.REQ.FUNC.EDITORIAL]}, \texttt{[TS.REQ.FUNC.PROB.BACKPACK.SI]},
\texttt{[TS.REQ.FUNC.PROB.BACKPACK.DEG]} (DOC-02) \\
\addlinespace
Приоритет & О --- обязательное \\
\addlinespace
Формулировка & Система должна формировать по запросу
\texttt{GET /api/tasks/backpack/editorial?type=\{тип\}} эталонный пример ранцевой
криптосистемы заданного типа (\texttt{BACKPACK\_CODE\_SUPER\_INCREASING} или
\texttt{BACKPACK\_CODE\_DEGREES}) с согласованным набором параметров: лёгкий рюкзак, модуль
$n$, множитель $\omega$, тяжёлый рюкзак, обратный множитель $\omega^{-1}$, шифртекст. \\
\addlinespace
Сценарий проверки & Для каждого допустимого типа проверить наличие всех параметров в ответе
и их математическую согласованность: лёгкий рюкзак сверхвозрастающий (для типа на степенях ---
элементы вида $p^i$); $n$ больше суммы элементов лёгкого рюкзака; $\omega$ взаимно просто с
$n$; $H_i = (L_i \cdot \omega) \bmod\ n$. Негативная ветка: запрос с недопустимым
\texttt{type} отклоняется. \\
\addlinespace
Зависимости & Требует аутентифицированного доступа --- SR-0XX (назначается в Спринте~2,
\texttt{[SR.AUTH]}). \\
\addlinespace
Владелец & \texttt{[SR.BACKPACK]}, Георгиев Д.Д. \\
\addlinespace
Статус & Черновик \\
\bottomrule
\end{tabularx}
\end{center}

\section{Функциональные требования по доменам \texttt{[SR.FUNC]}}\label{sec:func}
В настоящий раздел входят функциональные требования, сгруппированные по шести доменам
Системы; каждое требование оформляется карточкой по шаблону раздела~\ref{sec:card} и
снабжается маркером домена в заголовке соответствующего подраздела. Владельцы доменов
поставляют карточки для своих подразделов; редактор выполняет сквозную нумерацию и
единообразие (раздел~\ref{sec:own}).

\subsection{Аутентификация и пользователи \texttt{[SR.AUTH]}}\label{sec:auth}
Владелец: Кокорев В.Д. В подраздел войдут требования к: регистрации с валидацией тела
запроса и проверкой уникальности \texttt{email}; аутентификации по паре
\texttt{email}+пароль с выдачей JWT; проверке JWT на каждом запросе и ролям
(\texttt{ROLE\_USER}/\texttt{ADMIN}); профилю пользователя, привязке пользователя к учебной
группе и спискам групп; истории задач текущего пользователя.

\emph{Наполняется в Спринте~2 (тезисы --- в карточке домена \texttt{notes/domains/auth.md},
требования ТЗ --- \texttt{[TS.REQ.FUNC.AUTH.REG]}, \texttt{[TS.REQ.FUNC.AUTH.LOGIN]},
\texttt{[TS.REQ.FUNC.AUTH.JWT]}, \texttt{[TS.REQ.FUNC.GROUP.LIST]},
\texttt{[TS.REQ.FUNC.GROUP.JOIN]}, \texttt{[TS.REQ.FUNC.HISTORY]}).}
\begin{itemize}
  \item \texttt{SR-0XX} --- не назначено (регистрация, логин, JWT, роли, профиль, группы,
        история задач).
\end{itemize}

\subsection{Сессии задач и планировщик \texttt{[SR.SESSIONS]}}\label{sec:sessions}
Владелец: Анохин Г.А. В подраздел войдут требования к: созданию сессии задач по включённой
для группы студента конфигурации (не более одной активной сессии); генерации состава задач по
квалификаторам конфигурации; отправке и проверке решений с контролем владельца и состояния
задачи; переходам состояний задач \texttt{NEW}/\texttt{SOLVED}/\texttt{FAILED} и закрытию
сессии; автоматическому провалу просроченных задач планировщиком с интервалом 60~с при
лимите задачи 20~минут.

\emph{Наполняется в Спринте~2 (тезисы --- в карточке домена
\texttt{notes/domains/sessions.md}, требования ТЗ --- \texttt{[TS.REQ.FUNC.SESSION.CREATE]},
\texttt{[TS.REQ.FUNC.SESSION.TIME]}, \texttt{[TS.REQ.FUNC.SESSION.SUBMIT]}).}
\begin{itemize}
  \item \texttt{SR-0XX} --- не назначено (создание сессии, единственная активная сессия,
        отправка решения, переходы состояний, автопровал).
\end{itemize}

\subsection{Ранцевые криптосистемы \texttt{[SR.BACKPACK]}}\label{sec:backpack}
Владелец: Георгиев Д.Д. В подраздел войдут требования к: редакционным (эталонным) генерациям
двух типов (\texttt{GET/POST /api/tasks/backpack/editorial}); генерации задач ранцевой КС для
сессий (супервозрастающая последовательность, последовательность на степенях $p^i$, частично
скрытый лёгкий рюкзак); пошаговой проверке решений по контрольным соотношениям ранцевой КС.

\emph{Наполняется в Спринте~2 (тезисы --- в карточке домена
\texttt{notes/domains/backpack.md}, требования ТЗ --- \texttt{[TS.REQ.FUNC.PROB.BACKPACK.SI]},
\texttt{[TS.REQ.FUNC.PROB.BACKPACK.DEG]}).}
\begin{itemize}
  \item \texttt{SR-0XX} --- не назначено (редакционные генерации двух типов, генерация и
        проверка задач ранца в сессии).
\end{itemize}

\subsection{Бэкдор в ранцевой КС \texttt{[SR.BACKDOOR]}}\label{sec:backdoor}
Владелец: Минасов М.Р. В подраздел войдут требования к: генерации примера ранцевой КС на
степенях со встроенным бэкдором (\texttt{GET /api/backdoor/editorial}, параметры в теле
запроса) --- десятичные записи $\omega$ и $n$ встраиваются в младшие разряды элементов
тяжёлого рюкзака; ограничению числа попыток генерации; отсутствию сохранения результата в БД.

\emph{Наполняется в Спринте~2 (тезисы --- в карточке домена
\texttt{notes/domains/backdoor.md}, требования ТЗ --- \texttt{[TS.REQ.FUNC.PROB.BACKDOOR]}).}
\begin{itemize}
  \item \texttt{SR-0XX} --- не назначено (генерация примера с бэкдором, поведение при
        исчерпании попыток).
\end{itemize}

\subsection{Схема Эль-Гамаля \texttt{[SR.ELGAMAL]}}\label{sec:elgamal}
Владелец: Портнов Л.М. В подраздел войдут требования к: демонстрационному примеру с пошаговым
разбором (\texttt{GET /api/tasks/eigamal/editorial}); генерации учебного задания
(\texttt{GET /api/tasks/eigamal}); генерации и поэтапной проверке задач Эль-Гамаля в сессиях
(проверка $y$, $c_1$, $c_2$, расшифрования); математической библиотеке (генерация простых и
безопасных простых чисел, модулярная арифметика, поиск порождающего).

\emph{Наполняется в Спринте~2 (тезисы --- в карточке домена
\texttt{notes/domains/elgamal.md}, требования ТЗ --- \texttt{[TS.REQ.FUNC.PROB.ELGAMAL]}).}
\begin{itemize}
  \item \texttt{SR-0XX} --- не назначено (демо-пример, генерация задания, поэтапная проверка
        решения в сессии).
\end{itemize}

\subsection{Конфигурации задач и модерация \texttt{[SR.CONFIG]}}\label{sec:config}
Владелец: Самбьюи А.Ж. В подраздел войдут требования к: конфигурациям задач (имя и
упорядоченный список квалификаторов типов задач, задающий состав и порядок задач сессии);
включению конфигураций для групп (переключатель доступности); перечню конфигураций группы;
модерации --- просмотру сессий задач всех студентов; правилам доступа к функции модерации.
Требования по просмотру списка групп и вступлению в группы --- в
разделе~\ref{sec:auth} (\texttt{[SR.AUTH]}); здесь --- только конфигурации задач и их
привязка к группам.

\begin{sloppypar}
\emph{Наполняется в Спринте~2 (тезисы --- в карточке домена
\texttt{notes/domains/\allowbreak configurations.md}, требования ТЗ --- \texttt{[TS.REQ.FUNC.CONFIG.VIEW]}, \texttt{[TS.REQ.FUNC.MOD.VIEW]}).}
\end{sloppypar}
\begin{itemize}
  \item \texttt{SR-0XX} --- не назначено (состав конфигураций, включение группам, доступные
        конфигурации, модерация сессий).
\end{itemize}

\emph{Примечание. Требование \texttt{[TS.REQ.FUNC.EDITORIAL]} декомпозируется на
разделы~4.3--4.5 (учебные демо без записи в БД для каждого типа задач).}

\section{Нефункциональные требования \texttt{[SR.NFT]}}\label{sec:nft}
В настоящий раздел входят нефункциональные требования (НФТ) к Системе; формулировки НФТ
обязаны содержать измеримую меру (п.~\ref{sec:meth}). Карточка и правила заполнения --- те же,
что и для функциональных требований (раздел~\ref{sec:card}).

\subsection{Безопасность \texttt{[SR.SEC]}}\label{sec:sec}
\begin{sloppypar}
Владелец: Кокорев В.Д. В подраздел войдут требования к: хранению паролей (BCrypt); параметрам
JWT (секрет и срок жизни из переменных окружения, проверка на каждом запросе); разграничению
доступа по ролям, включая правила доступа к функции модерации; CORS; журналированию событий
безопасности; модели угроз с учётом известных особенностей реализации.
\end{sloppypar}

\begin{sloppypar}
\emph{Наполняется в Спринте~2 (тезисы --- карточки доменов \texttt{notes/domains/auth.md} и
\texttt{notes/domains/\allowbreak configurations.md}, требования ТЗ --- \texttt{[TS.REQ.SEC]}).}
\end{sloppypar}
\begin{itemize}
  \item \texttt{SR-0XX} --- не назначено (пароли, JWT, роли и модерация, CORS, модель угроз).
\end{itemize}

\subsection{Интеграция и интерфейсы \texttt{[SR.INT]}}\label{sec:int}
Владелец: Анохин Г.А. В подраздел войдут требования к: REST-интерфейсу Системы (базовый путь
\texttt{/api}, порт 8083), форматам запросов и ответов (JSON) и кодам состояния; переменным
окружения развёртывания; совместимости со СУБД PostgreSQL~17 и средой Docker; особым случаям
интерфейса (например, передача параметров в теле GET-запроса бэкдора).

\emph{Наполняется в Спринте~2 (тезисы --- карточки всех доменов, требования ТЗ ---
\texttt{[TS.REQ.INT]}).}
\begin{itemize}
  \item \texttt{SR-0XX} --- не назначено (REST-интерфейс, форматы и коды, переменные
        окружения, совместимость).
\end{itemize}

\subsection{Надёжность и эксплуатация \texttt{[SR.REL]}}\label{sec:rel}
Владелец: Самбьюи А.Ж. В подраздел войдут требования к: развёртыванию средствами
docker-compose (СУБД и приложение, контроль готовности БД); поведению при недоступности БД и
восстановлению после перезапуска; целостности схемы данных при обновлении; эксплуатационным
журналам; сохранности сессий задач при перезапуске планировщика автопровала.

\begin{sloppypar}
\emph{Наполняется в Спринте~2 (тезисы --- карточка домена
\texttt{notes/domains/\allowbreak configurations.md}, требования ТЗ --- \texttt{[TS.REQ.REL]}).}
\end{sloppypar}
\begin{itemize}
  \item \texttt{SR-0XX} --- не назначено (развёртывание, отказы и восстановление, журналы,
        сохранность данных).
\end{itemize}

\subsection{Производительность \texttt{[SR.PERF]}}\label{sec:perf}
Владелец: по решению редактора СТ (по плану Спринта~2 наполнение выполняет Самбьюи А.Ж.).
В подраздел войдут требования к: времени отклика эндпоинтов на типовой учебной нагрузке;
ресурсным характеристикам фонового планировщика автопровала; допустимым объёмам данных
сессий и истории задач.

\emph{Наполняется в Спринте~2 (требования ТЗ --- \texttt{[TS.REQ.PERF]}).}
\begin{itemize}
  \item \texttt{SR-0XX} --- не назначено (время отклика, нагрузка планировщика, объёмы
        данных).
\end{itemize}

\section{Трассировка \texttt{[SR.TRACE]}}\label{sec:trace}

\subsection{Правило инварианта трассировки \texttt{[SR.TRACE.RULE]}}
Инвариант комплекта: каждое SR трассируется минимум на одно SWR и минимум на один TC; каждое
SWR --- минимум на одно SR и минимум на один TC. Цепочка SR~$\to$~SWR~$\to$~TC ведётся в
матрице трассировки DOC-09 (CSV-источник, в комплект поставляется PDF) и проверяется на
каждой контрольной сверке и при приёмке 01.12.2026; нарушение инварианта --- блокирующее
замечание ФИ. TC создаётся только на уже поставленное SR/SWR; поле «Сценарий проверки»
карточки требования (раздел~\ref{sec:card}) является исходной основой соответствующего TC.

\subsection{Соответствие разделов СТ и разделов SWR \texttt{[SR.TRACE.MAP]}}
Детализация требований каждого домена выполняется в соответствующем документе программных
требований; фиксированное соответствие разделов приведено в таблице~\ref{tab:trace}. Матрица
DOC-09 заполняется инкрементально по мере готовности карточек (задача 2.14 Спринта~2).

\begin{table}[h!]
\centering
\caption{Соответствие разделов СТ, доменов-владельцев и разделов SWR}
\label{tab:trace}
\begin{tabularx}{\textwidth}{p{3.4cm}>{\raggedright\arraybackslash}p{6.6cm}>{\raggedright\arraybackslash}X}
\toprule
\textbf{Раздел СТ} & \textbf{Домен --- владелец} & \textbf{Раздел SWR} \\
\midrule
\texttt{[SR.AUTH]} & Аутентификация --- Кокорев В.Д. & DOC-06.1 \texttt{[SWR.AUTH]} \\
\texttt{[SR.SESSIONS]} & Сессии задач --- Анохин Г.А. & DOC-06.2 \texttt{[SWR.SESSIONS]} \\
\texttt{[SR.BACKPACK]} & Ранцевые КС --- Георгиев Д.Д. & DOC-06.3 \texttt{[SWR.BACKPACK]} \\
\texttt{[SR.BACKDOOR]} & Бэкдор --- Минасов М.Р. & DOC-06.4 \texttt{[SWR.BACKDOOR]} \\
\texttt{[SR.ELGAMAL]} & Эль-Гамаль --- Портнов Л.М. & DOC-06.5 \texttt{[SWR.ELGAMAL]} \\
\texttt{[SR.CONFIG]} & Конфигурации и модерация --- Самбьюи А.Ж. & DOC-06.6 \texttt{[SWR.CONFIG]} \\
\texttt{[SR.SEC]}, \texttt{[SR.INT]}, \texttt{[SR.REL]}, \texttt{[SR.PERF]} & НФТ --- по владельцам подразделов & DOC-06 сводно \texttt{[SWR.COMMON]} \\
\bottomrule
\end{tabularx}
\end{table}

\section{Порядок ведения документа}\label{sec:own}
\begin{enumerate}
  \item \textbf{Роли.} Редактор документа --- Анохин Г.А.: ведёт сквозную нумерацию и реестр
        SR, выполняет редакторскую сборку, контролирует соблюдение методики
        (раздел~\ref{sec:meth}) и единообразие карточек. Владельцы доменов
        (таблица~\ref{tab:markers}) поставляют карточки требований своих разделов по шаблону
        раздела~\ref{sec:card} и закрывают адресованные им замечания DEF.
  \item \textbf{Поставка карточек.} Карточки подаются владельцами в Спринте~2 (функциональные
        требования --- до 03.11.2026, НФТ --- до 04.11.2026) в тексте своих подразделов;
        редактор проверяет атомарность, трассируемость на \texttt{[TS.*]} и отсутствие дублей,
        после чего переводит статус карточек в «Согласовано».
  \item \textbf{Изменения.} Изменения базовой версии выполняются через лист регистрации
        изменений (титул, роли «Разработал/Согласовано/Утверждено») по процедурам ПУКПО
        (DOC-03). Контроль качества --- формальная инспекция ФИ-СТ: сессия 06.11.2026,
        протокол DOC-13.3 --- 07.11.2026; устранение замечаний ФИ --- первой задачей
        Спринта~3 (до 12.11.2026).
  \item \textbf{План версий} --- таблица~\ref{tab:versions}.
\end{enumerate}

\begin{table}[h!]
\centering
\caption{План версий документа}
\label{tab:versions}
\begin{tabularx}{\textwidth}{p{1.6cm}p{2.4cm}>{\raggedright\arraybackslash}X}
\toprule
\textbf{Версия} & \textbf{Дата} & \textbf{Содержание} \\
\midrule
v0.1 & 24.10.2026 & Скелет документа (Спринт~1): структура, методика ID, шаблон карточки
требования, примеры-образцы SR-001--SR-003. \\
v1.0 & 05.11.2026 & Редакторская сборка после наполнения: функциональные требования шести
доменов, НФТ, сквозная нумерация и реестр SR (Спринт~2). \\
v1.1 & 12.11.2026 & Устранение замечаний ФИ-СТ (DOC-13.3), закрытие DEF. \\
\bottomrule
\end{tabularx}
\end{table}

\end{document}
